\chapter{Operator, Kernel, and Spectral Realization Programme}
\label{ch:operator}

\section{Hilbert--Pólya criterion}
A standard sufficient route to the Riemann hypothesis is to construct a self-adjoint operator $\widehat H$ whose spectrum consists of the ordinates of the nontrivial zeros. In its cleanest form, one seeks a relation such as
\begin{equation}
D(E)=C(E)\,\xi\!\left(\frac12+iE\right),
\label{eq:spectral-det}
\end{equation}
where $D(E)$ is a spectral determinant of $\widehat H$ and $C(E)$ is nonzero in the relevant domain. If $\widehat H$ is self-adjoint, its eigenvalues $E$ are real. The corresponding zeros then have real part $1/2$.

The principal difficulty is not writing an operator with a spectrum fitted to a pre-existing list of zeros. It is constructing a natural operator whose domain, boundary conditions, self-adjointness, and determinant are derived independently from arithmetic or geometry. Berry and Keating related the smooth zero-counting function to the semiclassical Hamiltonian $H=xp$, while Connes developed an absorption-spectrum interpretation through noncommutative geometry \cite{BerryKeating1999,Connes1999}. These are structural precedents rather than completed proofs.

\section{Information-spinor operator ansatz}
The present framework suggests a tensor-product Hilbert space
\begin{equation}
\mathcal H=\mathcal H_{\mathrm{arith}}\otimes\mathbb C^2,
\end{equation}
where the two-dimensional factor carries the complement degree of freedom. Let $\tau_x,\tau_y,\tau_z$ denote Pauli matrices acting on the complement factor. A schematic operator may be written
\begin{equation}
\widehat{\mathcal D}
=\widehat H_{\mathrm{arith}}\otimes I
+\widehat V\otimes\tau_x
+\widehat W\otimes\tau_y
+\widehat M\otimes\tau_z.
\label{eq:dirac-ansatz}
\end{equation}
The exchange involution acts as $\tau_x$ combined with complex conjugation. The balance axis corresponds to vanishing expectation of $\tau_z$:
\begin{equation}
\langle\tau_z\rangle=2\sigma-1=0.
\end{equation}
A chiral or particle-hole type symmetry could force eigenstates at a zero to have equal complement populations.

This is only an architecture. To become useful, the arithmetic operators must be specified and the determinant or scattering matrix must be shown to reproduce $\xi$.

\section{Positive balance penalty}
The Shannon deficit defines a nonnegative scalar potential
\begin{equation}
V_S(\sigma)=\ln2-\Hbin(\sigma).
\end{equation}
Likewise, the phase-$\pi$ cancellation residual is
\begin{equation}
V_I(\sigma)=\left(\sqrt\sigma-\sqrt{1-\sigma}\right)^2.
\end{equation}
Both vanish uniquely at $1/2$ and agree to quadratic order. An effective transverse Hamiltonian could contain
\begin{equation}
\widehat H_\perp
=-\alpha\frac{\dd^2}{\dd\sigma^2}
+\beta V_S(\sigma)
+\gamma V_I(\sigma),
\end{equation}
with suitable boundary conditions on $(0,1)$. Its ground state would be centered at $1/2$, but this alone would not prove that all zeta spectral states have zero transverse displacement. Excited transverse states generally exist. A supersymmetry, projection, topological index, or constraint would be needed to restrict the physical sector to the balanced mode.

\section{Holonomy as a boundary condition}
The Berry connection of Chapter \ref{ch:bloch} suggests imposing a monodromy condition
\begin{equation}
\Psi(\phi+2\pi)=-\Psi(\phi).
\label{eq:antiperiodic}
\end{equation}
Antiperiodic boundary conditions yield half-integer Fourier modes
\begin{equation}
\Psi_m(\phi)=e^{i(m+1/2)\phi},
\qquad m\in\mathbb Z.
\end{equation}
This is a precise mechanism by which a half shift appears spectrally. The challenge is to derive the zeta real part from this half-shift rather than merely reproducing a familiar spin structure.

One possible target is a transfer operator or Dirac-type operator whose Mellin exponent is shifted by the spin structure:
\begin{equation}
s=\frac12+iE.
\end{equation}
If the $1/2$ is a fixed spin-connection contribution and $E$ is the real spectrum of a self-adjoint generator, RH follows. The unresolved tasks are to identify the arithmetic evolution and prove the determinant identity \eqref{eq:spectral-det}.

\section{Mellin generator and dilation}
The classical generator of dilations is $xp$. Quantum mechanically, a symmetric realization is
\begin{equation}
\widehat D=\frac12(\widehat x\widehat p+\widehat p\widehat x)
=-i\left(x\frac{\dd}{\dd x}+\frac12\right)
\end{equation}
on an appropriate domain. The appearance of $1/2$ here is forced by symmetric ordering and by unitarity of dilations in $L^2(\mathbb R_+,\dd x)$. Indeed, the unitary dilation group acts as
\begin{equation}
(U_a f)(x)=e^{a/2}f(e^a x),
\end{equation}
and its generator contains the half-density term.

This provides a second exact origin of the half shift, independent of Shannon entropy. Mellin eigenfunctions behave as
\begin{equation}
f_E(x)=x^{-1/2+iE},
\end{equation}
up to conventions. Their exponent lies on the critical line. Berry--Keating type models exploit this dilation structure \cite{BerryKeating1999,Sierra2007}.

\begin{proposition}[Half-density origin of the critical line]
The unitary representation of dilations on $L^2(\mathbb R_+,\dd x)$ has generalized eigenfunctions with Mellin exponent $-1/2+iE$. Under the conventional zeta Mellin variable, this corresponds to a fixed real part $1/2$.
\end{proposition}

The proposition explains why self-adjoint dilation generators naturally live on the critical line. It still does not establish that their spectral boundary condition is the zeta function.

\section{Kernel criterion}
Another route starts with the cosine transform \eqref{eq:xi-fourier}. A sufficient result would be a theorem placing the kernel $\Phi$ in a class whose Fourier transform has only real zeros. Candidate concepts include Pólya frequency functions, total positivity, variation-diminishing transforms, and de Branges spaces.

The information interpretation suggests examining whether $\Phi$ can be normalized as a probability or amplitude kernel whose complement operation has a unique fixed point. The relevant test is not whether $\Phi$ is positive; positivity alone is insufficient to force real zeros of a Fourier transform. One needs a much stronger structural property.

\section{De Branges strategy}
De Branges theory gives conditions under which an entire function belongs to a Hermite--Biehler class and has controlled zero geometry. A possible target is to construct an entire function $E(z)$ such that
\begin{equation}
\Xi(z)=\frac12\left(E(z)+E^{\#}(z)\right),
\qquad
E^{\#}(z)=\overline{E(\overline z)},
\end{equation}
with
\begin{equation}
|E(z)|>|E^{\#}(z)|
\qquad \text{for }\Imm z>0.
\end{equation}
Such an inequality would encode a positivity or causality condition. In the two-branch language, $E$ and $E^{\#}$ are complementary amplitudes, and the real axis is their fixed boundary. Establishing the Hermite--Biehler inequality for an appropriate zeta-derived $E$ would be a substantive proof route rather than a restatement of symmetry.

\section{Minimal proof obligations}
A publishable operator claim must specify:
\begin{enumerate}
\item the Hilbert space and inner product;
\item the dense operator domain;
\item the adjoint and deficiency indices;
\item the self-adjoint extension, if any;
\item the boundary conditions and why they are arithmetic rather than fitted;
\item the determinant or scattering function;
\item proof that any prefactor multiplying $\xi$ is nonzero in the critical strip;
\item proof that no tabulated zero data were inserted into the operator.
\end{enumerate}
Without these items, an operator analogy remains heuristic.

\chapter{Correlation-kernel density criterion}\label{ch:correlation-kernel}

\section{A second analytic route from the Xi kernel}

The canonical Fourier-kernel representation introduced in the operator programme admits a second route that is independent of the Hermite--Biehler candidate audited in XF-2.  For the even Riemann kernel $\Phi$ and
\[
\Xi(z)=\int_{\mathbb R}\Phi(t)e^{izt}\,\dd t,
\]
define the second correlation function
\begin{equation}\label{eq:nu2-correlation}
\nu_2(t)
=
\int_{-\infty}^{\infty}
(t-2s)^2\Phi(t-s)\Phi(s)\,\dd s.
\end{equation}
For $0<|y|<1/2$, set
\begin{equation}\label{eq:phi2y-correlation}
\Phi_{2,y}(t)=\cosh(ty)\,\nu_2(t).
\end{equation}
These are the $n=2$ correlation objects used by Dimitrov and Xu in their Fourier--Wronskian criteria \cite{DimitrovXu2016}.

\section{Published RH-equivalent density criterion}

Dimitrov--Xu Theorem~1.1 supplies a standard external equivalence.  In the present notation,
\begin{equation}\label{eq:dx-density-rh}
\boxed{
\mathrm{RH}
\iff
\forall\,y\in\left(-\frac12,\frac12\right)\setminus\{0\},\quad
\overline{\operatorname{span}\mathcal T(\Phi_{2,y})}^{\,L^1(\mathbb R)}
=L^1(\mathbb R),
}
\end{equation}
where $\mathcal T(f)$ denotes the family of all translates of $f$.

The same work supplies an equivalent convolution-annihilator formulation: for every admissible $y$, the condition
\begin{equation}\label{eq:dx-annihilator}
\Phi_{2,y}*g=0,\qquad g\in L^\infty(\mathbb R),
\end{equation}
forces $g=0$ exactly on the RH surface \cite{DimitrovXu2016}.

Within the typed implication graph, the theorem itself has status \status{STANDARD}.  The global density condition and the equivalent annihilator condition have status \status{OPEN}; they are the mathematical premises that must be independently established before the standard equivalence can close the RH node.

\section{Wronskian bridge}

For $n=2$, the Fourier--Wronskian identity gives
\begin{equation}\label{eq:dx-wronskian}
W_2(\Xi;x)
=
\Xi(x)\Xi''(x)-\Xi'(x)^2
=
-\mathcal F[\nu_2](x),
\end{equation}
with the Fourier convention of \cite{DimitrovXu2016}.  Since $\Xi$ is real on the real $z$ axis, its first Laguerre quantity obeys
\begin{equation}\label{eq:laguerre-wronskian}
\boxed{
L_1[\Xi](x)
=
\Xi'(x)^2-\Xi(x)\Xi''(x)
=
-W_2(\Xi;x)
=
\mathcal F[\nu_2](x).
}
\end{equation}
This gives an exact consistency bridge between the correlation-kernel and derivative representations.

\section{Executable diagnostic layer}

The implementation in \texttt{src/critical\_axis/correlation\_kernel.py} evaluates $\nu_2$, $\Phi_{2,y}$, the Wronskian, and the Laguerre quantity with explicit finite truncation controls.  The validation suite checks:
\begin{enumerate}[label=(\roman*)]
\item evenness and positive reference values of the implemented kernel;
\item numerical evenness of $\nu_2$;
\item the exact $y\leftrightarrow-y$ symmetry of $\Phi_{2,y}$;
\item real-axis closure of Eq.~\eqref{eq:laguerre-wronskian};
\item fail-closed enforcement of $0<|y|<1/2$;
\item sampled Laguerre positivity as a \status{NUMERICAL\_DIAGNOSTIC} quantity.
\end{enumerate}

Finite samples and finite translate families carry diagnostic authority only.  The repository records this layer as
\[
\status{NUMERICAL\_DIAGNOSTIC}.
\]
The theorem-level target retains its full quantifiers over every admissible $y$, all translates, and the complete space $L^1(\mathbb R)$.

\section{Relation to the branch-cancellation programme}

The two active routes now share the same canonical Xi kernel but expose different closure obligations:
\begin{align}
\Phi
&\longrightarrow A_\pm
\longrightarrow \text{exact zero cancellation}
\longrightarrow \text{population/coordinate bridge},\\
\Phi
&\longrightarrow \nu_2
\longrightarrow \Phi_{2,y}
\longrightarrow \text{$L^1$ translation density or annihilator closure}.
\end{align}
The first route isolates a local zero-state representation problem.  The second route converts RH into a global harmonic-analysis completeness problem.  Their coexistence supplies a useful cross-audit: a future mechanism that closes the critical axis should be checked against both the local branch structure and the global correlation-kernel criterion.

\section{Current frontier}

The immediate XF-3 proof obligation is therefore
\begin{equation}\label{eq:xf3-frontier}
\boxed{
\forall\,0<|y|<\frac12:\quad
\overline{\operatorname{span}\mathcal T(\Phi_{2,y})}^{\,L^1(\mathbb R)}
=L^1(\mathbb R),
}
\end{equation}
or equivalently the global bounded-annihilator condition of Eq.~\eqref{eq:dx-annihilator}.  The solver records both as \status{OPEN\_RH\_EQUIVALENT\_CRITERION}; the Riemann-hypothesis node therefore retains status \status{OPEN}.

\chapter{Wiener--Laguerre scalar criterion}\label{ch:wiener-laguerre-scalar}

\section{Reduction of the density frontier}

Chapter~\ref{ch:correlation-kernel} expresses the Dimitrov--Xu critical-axis criterion as density of the translates of
\[
\Phi_{2,y}(t)=\cosh(ty)\nu_2(t),
\qquad
0<|y|<\frac12.
\]
Wiener's $L^1$ translation theorem allows this infinite-dimensional statement to be reduced to a pointwise condition on a Fourier transform.  The Wronskian identities of Dimitrov and Xu then identify that transform directly with derivatives of the completed Xi function \cite{DimitrovXu2016}.

For real $x$ and admissible $y$, define
\begin{equation}\label{eq:q-xi-definition}
Q_\Xi(x,y)
:=
|\Xi'(x+iy)|^2
-\Ree\!\left(
\Xi(x+iy)\overline{\Xi''(x+iy)}
\right).
\end{equation}

\section{Wronskian identity}

Write
\[
\Xi(x+iy)=\phi(x,y)+i\psi(x,y).
\]
The Cauchy--Riemann equations give
\begin{align}
Q_\Xi(x,y)
={}&\phi_x(x,y)^2-\phi(x,y)\phi_{xx}(x,y)\\
&+\psi_x(x,y)^2-\psi(x,y)\psi_{xx}(x,y).
\end{align}
With $W_2(f)=ff''-(f')^2$, this becomes
\begin{equation}\label{eq:q-xi-wronskian}
Q_\Xi(x,y)
=-W_2(\phi(\cdot,y);x)-W_2(\psi(\cdot,y);x).
\end{equation}
Dimitrov--Xu Theorem~1.3 / Theorem~2.5 and their Section~3 correlation identity yield
\begin{equation}\label{eq:q-xi-fourier}
\boxed{
Q_\Xi(x,y)=\mathcal F[\Phi_{2,y}](x).
}
\end{equation}
The repository records Eq.~\eqref{eq:q-xi-fourier} as \status{STANDARD\_PASS}.

\section{Tauberian equivalence}

For each fixed admissible $y$, Wiener's $L^1$ theorem states
\begin{equation}\label{eq:wiener-density-nonzero}
\overline{\operatorname{span}\mathcal T(\Phi_{2,y})}^{L^1(\mathbb R)}
=L^1(\mathbb R)
\iff
\mathcal F[\Phi_{2,y}](x)\neq0
\quad\forall x\in\mathbb R.
\end{equation}
The Dimitrov--Xu proof also gives strict positivity at the origin,
\begin{equation}\label{eq:q-origin-positive}
Q_\Xi(0,y)>0
\qquad
\left(0<|y|<\frac12\right).
\end{equation}
Because $Q_\Xi(\cdot,y)$ is continuous and real-valued on $\mathbb R$, a globally nonvanishing $Q_\Xi$ has the positive sign selected by Eq.~\eqref{eq:q-origin-positive}.  Therefore
\begin{equation}\label{eq:density-q-positive}
\boxed{
\overline{\operatorname{span}\mathcal T(\Phi_{2,y})}^{L^1}=L^1
\iff
Q_\Xi(x,y)>0\quad\forall x\in\mathbb R.
}
\end{equation}
Combining Eq.~\eqref{eq:density-q-positive} with Dimitrov--Xu Theorem~1.1 gives
\begin{theorem}[Wiener--Laguerre scalar RH criterion]\label{thm:wiener-laguerre-rh}
Under the standard Xi-kernel hypotheses used by Dimitrov and Xu,
\begin{equation}\label{eq:q-rh-equivalence}
\boxed{
\mathrm{RH}
\iff
\forall\,0<|y|<\frac12\ \forall x\in\mathbb R:\quad
Q_\Xi(x,y)>0.
}
\end{equation}
\end{theorem}

The equivalence in Theorem~\ref{thm:wiener-laguerre-rh} has status \status{STANDARD}.  The globally quantified sign condition has status
\[
\status{OPEN\_RH\_EQUIVALENT\_CRITERION}.
\]

\section{Executable diagnostic}

The function
\begin{center}
\texttt{xi\_wiener\_laguerre\_scalar(x, y)}
\end{center}
implements Eq.~\eqref{eq:q-xi-definition} and enforces the open domain $0<|y|<1/2$.  The tests verify $y\leftrightarrow-y$ symmetry, positive origin reference values, and fail-closed domain handling.

Finite samples retain \status{NUMERICAL\_DIAGNOSTIC} authority.  An analytic closure must carry both universal quantifiers in Eq.~\eqref{eq:q-rh-equivalence}.

\section{Certified falsification surface}

The scalar form provides a direct falsification target.  A certified pair $(x_*,y_*)$ with
\[
x_*\in\mathbb R,
\qquad
0<|y_*|<\frac12,
\qquad
Q_\Xi(x_*,y_*)\le0
\]
combined with Eq.~\eqref{eq:q-origin-positive} and continuity would force a real zero of $Q_\Xi(\cdot,y_*)=\mathcal F[\Phi_{2,y_*}]$.  The corresponding Wiener density condition would fail at that $y_*$.

This makes the next research frontier sharply local in representation while globally quantified in domain:
\begin{equation}\label{eq:xf4-frontier}
\boxed{
Q_\Xi(x,y)>0
\quad
\text{for every }x\in\mathbb R,
\quad0<|y|<\frac12.
}
\end{equation}
The next analytic task is to seek a structural representation of $Q_\Xi$ that certifies this sign over the entire strip.

\chapter{Nonlocal transverse curvature}\label{ch:nonlocal-curvature}

\section{From the Wiener--Laguerre scalar to geometry}

Chapter~\ref{ch:wiener-laguerre-scalar} isolates the RH-equivalent scalar condition
\begin{equation}\label{eq:xf5-q-frontier}
Q_\Xi(x,y)>0,
\qquad
x\in\mathbb R,
\qquad
0<|y|<\frac12.
\end{equation}
The next step is an exact geometric reformulation of the same scalar.  Set
\[
F_x(y):=|\Xi(x+iy)|^2.
\]
Since $\Xi$ is real-entire,
\[
\overline{\Xi(x+iy)}=\Xi(x-iy),
\]
and direct differentiation yields
\begin{equation}\label{eq:xf5-curvature-identity}
\boxed{
\frac12 F_x''(y)
=
|\Xi'(x+iy)|^2
-\Ree\!\left(
\Xi(x+iy)\overline{\Xi''(x+iy)}
\right)
=Q_\Xi(x,y).
}
\end{equation}
Equation~\eqref{eq:xf5-curvature-identity} has status \status{EXACT}.  Therefore the XF-4 global sign frontier is equivalent to strict transverse convexity,
\begin{equation}\label{eq:xf5-convexity-frontier}
\boxed{
F_x''(y)>0
\quad
\text{for every }x\in\mathbb R,
\quad0<|y|<\frac12.
}
\end{equation}
The universal statement in Eq.~\eqref{eq:xf5-convexity-frontier} retains
\[
\status{OPEN\_RH\_EQUIVALENT\_CRITERION}.
\]

\section{Connection with the classical growth criterion}

The function $F_x$ is even in $y$, hence
\[
F_x'(0)=0.
\]
Strict convexity on $0<y<1/2$ therefore implies
\begin{equation}\label{eq:xf5-inner-growth}
F_x'(y)>0,
\qquad 0<y<\frac12.
\end{equation}
The outer region $y\ge 1/2$, corresponding to $\Ree s\ge1$, is controlled by the standard symmetric Hadamard/logarithmic-derivative argument.  Combining the inner and outer regions recovers the classical positivity/growth route associated with Hinkkanen and Lagarias \cite{Lagarias1999,Planat2026Theta}.

The solver records the same implication chain in four breakable stages:
\begin{enumerate}
\item \path{Q_Xi strict positivity};
\item equivalently, strict transverse convexity;
\item consequently, positive vertical growth;
\item with the outer-halfplane closure, the classical RH-equivalent growth criterion.
\end{enumerate}
The first globally quantified condition retains \status{OPEN}.  Thus XF-5 changes the representation of the proof obligation and preserves its evidential status.

\section{Differentiated theta-kernel representation}

The repository normalization is
\begin{equation}\label{eq:xf5-repo-fourier}
\Xi(z)=2\int_0^\infty\Phi(u)\cos(zu)\,\dd u.
\end{equation}
Define
\[
D(z)=\int_0^\infty\Phi(u)\cos(zu)\,\dd u=\frac12\Xi(z).
\]
In diagonal coordinates
\[
u=a+b,
\qquad
v=a-b,
\qquad
a>|b|,
\]
write
\[
M(a,b)=\Phi(a+b)\Phi(a-b).
\]
The symmetrized first-growth kernel can be represented as
\begin{equation}\label{eq:xf5-growth-kernel}
K_{x,y}(a,b)
=
\frac a2\cos(2xb)\sinh(2ya)
+
\frac b2\cos(2xa)\sinh(2yb).
\end{equation}
Differentiation gives the local curvature kernel
\begin{equation}\label{eq:xf5-curvature-kernel}
\boxed{
L_{x,y}(a,b)
:=\partial_y K_{x,y}(a,b)
=
a^2\cos(2xb)\cosh(2ya)
+
b^2\cos(2xa)\cosh(2yb).
}
\end{equation}
Under Eq.~\eqref{eq:xf5-repo-fourier}, the XF-4 scalar becomes
\begin{equation}\label{eq:xf5-nonlocal-q}
\boxed{
Q_\Xi(x,y)
=4\iint_{a>|b|}
M(a,b)L_{x,y}(a,b)\,\dd a\,\dd b.
}
\end{equation}
The differentiated-kernel identity and normalization bookkeeping have status \status{EXACT}.  The sign of the complete integral over the full admissible domain remains the active global obligation.

\section{Exact positive anchor at the central slice}

For $x=0$,
\begin{equation}\label{eq:xf5-xzero-kernel}
L_{0,y}(a,b)
=a^2\cosh(2ya)+b^2\cosh(2yb)>0
\end{equation}
throughout the interior domain $a>|b|$, $a>0$.  Together with the positive theta mass $M(a,b)$, this gives pointwise positivity of the XF-5 curvature integrand at the central slice.  It supplies a kernel-level version of the positive-sign anchor used in the Wiener argument of Chapter~\ref{ch:wiener-laguerre-scalar}.

\section{Nonlocality firewall}

Planat's 2026 theta-kernel analysis derives an exact longitudinal--transverse decomposition of the Xi growth derivative \cite{Planat2026Theta}.  The companion August 2026 preprint proves, for that phase-aligned decomposition, that the two sectors cancel to all algebraic orders and that universal independent block positivity fails: sufficiently far in the oscillatory regime at least one aligned block is negative \cite{Planat2026Nonlocal}.

The repository therefore records the following research firewall:
\begin{itemize}
\item Global correlated sign control --- \status{ACTIVE}.
\item Exact resummation preserving sector correlation --- \status{ACTIVE}.
\item Positive representation of the fully coupled integral --- \status{ACTIVE}.
\item Universal independent phase-aligned block positivity.\par\noindent Status: \status{LITERATURE\_NO\_GO}.
\item Finite block sampling --- \status{NUMERICAL\_DIAGNOSTIC}.
\end{itemize}
The no-go result is representation-specific.  It narrows the admissible strategy to globally correlated control of Eq.~\eqref{eq:xf5-nonlocal-q} or to a different exact representation that preserves the sign-carrying remainder.

\section{Executable audit surface}

The module \path{critical_axis.nonlocal_curvature} exposes
\begin{itemize}
\item \path{xi_transverse_curvature(x,y)}, implementing Eq.~\eqref{eq:xf5-curvature-identity};
\item \path{xi_transverse_curvature_direct(x,y)}, independently differentiating $|\Xi|^2$ for regression;
\item \path{theta_growth_kernel(x,y,a,b)};
\item \path{theta_curvature_kernel(x,y,a,b)}, implementing Eq.~\eqref{eq:xf5-curvature-kernel};
\item \path{theta_curvature_integrand}, evaluating one local contribution with fail-closed diagonal coordinates.
\end{itemize}

The corresponding tests lock the factor of two in Eq.~\eqref{eq:xf5-curvature-identity}, compare analytic and direct second derivatives, differentiate Eq.~\eqref{eq:xf5-growth-kernel} independently, verify the $x=0$ positive anchor, and reject invalid $a,b$ domains.

\section{XF-5 frontier}

The new proof surface is therefore the complete correlated curvature integral
\begin{equation}\label{eq:xf5-final-frontier}
\boxed{
\iint_{a>|b|}
M(a,b)L_{x,y}(a,b)\,\dd a\,\dd b>0
}
\end{equation}
for every real $x$ and every $0<|y|<1/2$.  Viable next steps are an analytic global lower bound, an exact correlation-preserving resummation, a positive-definite representation of the fully coupled integral, or a certified counterexample on the admissible domain.  The universal sign statement retains \status{OPEN\_RH\_EQUIVALENT\_CRITERION}.

